\documentclass[11pt]{article}

\usepackage[margin=1.15in]{geometry}
\usepackage{amsmath,amssymb}
\usepackage{booktabs}
\usepackage[numbers,sort&compress]{natbib}
\usepackage{xcolor}
\usepackage[colorlinks=true,linkcolor=blue!60!black,citecolor=blue!60!black,urlcolor=blue!60!black]{hyperref}
\usepackage{microtype}

\newcommand{\br}[1]{\langle #1 \rangle}

\title{Assembly of the two-loop MHV integrand from maximally positive
mutual numerators at five and six points}

\author{Jon Torrez\thanks{Correspondence: \texttt{epg@subnet345.com}.
Computations performed by an AI agent swarm---a mix of OpenAI and Anthropic
models---working on quantum physics under the author's direction, with
results independently cross-checked between agents; all code at
\url{https://github.com/jrtorrez31337/executable-positive-geometry}.}}

\date{July 5, 2026}

\begin{document}
\maketitle

\begin{abstract}
For the two-loop four-point MHV amplituhedron, the canonical form on
each positive chart is known to be a pure dlog form times the ratio
$D^{+}/D$, where $D = \br{ABCD}$ is the mutual positivity condition and
$D^{+}$ is the sum of its positive monomials
\citep{Kojima:2020gxs}. We test, by exact computation, whether this
``maximally positive numerator'' structure \emph{assembles} at higher
multiplicity: whether the full two-loop MHV integrand at $n=5,6$ equals
the sum over all composite cells (ordered pairs of one-loop sign-flip
cells) of Kermit $\times$ Kermit $\times\, D^{+}/D$, with no
corrections. It does. At $n=5$ (9 composite cells) and $n=6$ (36
composite cells) the assembly reproduces the local double-pentagon
representation of the integrand \citep{ArkaniHamed:2010gh} with ratio
exactly $1$ in rational arithmetic at independent generic kinematic
points; interior seams cancel pairwise; and all well-defined transverse
octa-cut residues at $n=5$ are quantized as required. Two structural
observations explain why the preconditions cannot fail at any $n$: in
the sign-flip charts, $\br{ABCD}$ is automatically multilinear in the
eight chart variables, and every monomial coefficient is a single
Pl\"ucker bracket, so the sign split defining $D^{+}$ is independent of
the choice of positive kinematics. In particular the quasi-linear
complication anticipated in \citep{Kojima:2020gxs} does not arise for
two-loop MHV in these raw charts at $n \le 6$, suggesting the naive
assembly may hold for all $n$. All computations are exact and
reproducible from the public repository.
\end{abstract}

\section{Setup}

We use momentum twistors $Z_1,\dots,Z_n$ on a moment curve (all ordered
Pl\"ucker brackets positive) and the one-loop MHV amplituhedron in its
sign-flip description \citep{ArkaniHamed:2013jha,ArkaniHamed:2013kca,
Kojima:2020gxs}: a loop line $AB$ with $\br{AB\,i\,i{+}1} > 0$ (twisted
cyclic) and exactly two sign flips in the sequence $\br{AB1i}$,
$i = 2,\dots,n$. The flip positions $(i,j)$, $2 \le i < j \le n-1$,
label $(n-2)(n-3)/2$ cells, each with the positive chart
\begin{equation}
A = Z_1 + x\,Z_i + w\,Z_{i+1}, \qquad B = -Z_1 + y\,Z_j + z\,Z_{j+1},
\qquad x,w,y,z > 0 .
\end{equation}
On each cell the canonical form is the corresponding ``Kermit'' term
\citep{ArkaniHamed:2010gh},
\begin{equation}
K_{ij}(A,B) \;=\;
\frac{\br{AB\,(1\,i\,i{+}1)\cap(1\,j\,j{+}1)}^{2}}
{\br{AB1i}\br{AB\,i\,i{+}1}\br{AB\,1\,i{+}1}\,
 \br{AB1j}\br{AB\,j\,j{+}1}\br{AB\,1\,j{+}1}},
\end{equation}
which we verify to be an exact dlog form on its chart
($K_{ij}\,\br{AB\,i\,i{+}1}\br{AB\,j\,j{+}1}\,xwyz = 1$ identically)
for all cells at $n = 5, 6$.

At two loops, the geometry is the set of ordered pairs of lines
$(AB, CD)$, each in the one-loop region, with the single mutual
condition $\br{ABCD} > 0$ \citep{ArkaniHamed:2013kca,Kojima:2020gxs,
Dian:2024tvi}. Composite cells are ordered pairs of one-loop cells;
in each composite chart write $D = \br{ABCD}$ as a polynomial in the
eight chart variables and let $D^{+}$ denote the sum of its
positive-coefficient monomials.

\paragraph{Two structural facts.} (i) $D$ is automatically multilinear:
each chart variable enters exactly one column of the $4\times4$
determinant. (ii) Each monomial coefficient is a single Pl\"ucker
bracket $\pm\br{Z_aZ_bZ_cZ_d}$ (one basis vector chosen per column), so
its sign is fixed by the ordering of the moment curve alone: the split
$D = D^{+} - D^{-}$ is independent of the choice of positive $Z$'s. We
verified both facts computationally for all $9$ composite cells at
$n=5$ and all $36$ at $n=6$ (monomial counts ranging from 8 to 26), in
addition to their one-line proofs. Consequently the ``maximally
positive numerator'' is well defined at every $n$, and the only
question is whether the assembly reproduces the integrand.

\section{The assembly and its tests}

\paragraph{Postulate.} With all quantities evaluated at common
representatives of the loop lines,
\begin{equation}
\label{eq:P1}
M_2^{(n)} \;=\; \sum_{\text{composite cells}}
K_{c_1}(A,B)\; K_{c_2}(C,D)\; \frac{D^{+}_{c_1c_2}}{D_{c_1c_2}},
\end{equation}
where $D_{c_1c_2}$, $D^{+}_{c_1c_2}$ are evaluated by canonically
extracting each pair's chart coordinates from the lines, and the sum is
over \emph{ordered} pairs (so $M_2 = 2\,\mathcal{A}^{2\text{-loop}}$ in
the convention where the physical amplitude carries $1/2$ for
indistinguishable loops). At $n=4$ this is the known result
\citep{Kojima:2020gxs}, which our implementation reproduces exactly.

\paragraph{Gold standard.} We compare \eqref{eq:P1} against the local
double-pentagon representation \citep{ArkaniHamed:2010gh},
$\mathcal{A}^{2\text{-loop}}_{\rm MHV} = \tfrac12\sum_{i<j<k<l<i} I_{ijkl}$,
with the chiral numerators
$\br{AB\,(i{-}1\,i\,i{+}1)\cap(j{-}1\,j\,j{+}1)}\br{ijkl}
\br{CD\,(k{-}1\,k\,k{+}1)\cap(l{-}1\,l\,l{+}1)}$, implemented verbatim
from the reference. In exact rational arithmetic at independent generic
kinematic points:

\begin{center}
\begin{tabular}{lccc}
\toprule
$n$ & composite cells & double pentagons & ratio
$M_2 / \sum I$ \\
\midrule
$5$ & $9$  & $20$ & $1$ (three points) \\
$6$ & $36$ & $60$ & $1$ (two points) \\
\bottomrule
\end{tabular}
\end{center}

The ratios are exact rational numbers, not floating-point agreements.

\paragraph{Seams.} Individual composite-cell terms have poles on the
spurious hyperplanes $\br{AB1i}$ ($i$ non-adjacent). Along exact
one-parameter rational paths onto each seam at $n=5$, individual terms
diverge as $1/\varepsilon$ while the assembled sum is constant to six
digits across four orders of magnitude in $\varepsilon$, and the sum
diverges as required on the physical mutual pole $\br{ABCD}\to 0$.

\paragraph{Leading singularities ($n=5$).} On the two standard
octa-cut families with rational Schubert solutions, all well-defined
transverse residues of the assembly are exactly $0$, in the two ways
the physics demands: $50/50$ zeros on kissing-pentagon cuts (at $n=5$
all disjoint non-adjacent chords cross, so every transverse kissing cut
is non-planar), and $20/20$ zeros on the plane-intersection branches of
pentagon-box cuts (where the chiral numerators are designed to vanish).
The remaining chord-branch pentagon-box configurations at $n=5$ are
composite strata (cut solutions pass through polygon vertices; naive
transverse residues are ray-dependent, which our procedure detects
automatically); their iterated-residue census is left open. The residue
machinery is calibrated in-code on the one-loop box quadruple cut
(residues exactly $\mp 1$, invariant under representative changes).

\section{Discussion}

Three comments. First, the result is a verification with an unexpected
edge: the quasi-linear machinery of \citep{Kojima:2020gxs}, developed
because ``the tedious triangulation is inevitable'' in general, is
never needed for two-loop MHV in the raw sign-flip charts at
$n \le 6$ --- the naive term-by-term positive part assembles into the
exact integrand. Given the structural facts (i)--(ii), which hold at
every $n$, it is natural to conjecture that \eqref{eq:P1} holds for all
$n$ at two loops MHV; we have no proof, and the composite
leading-singularity sector remains unchecked beyond the families
listed. Second, two predictions made during this work were refuted by
its own tests and are recorded in the repository alongside the
confirmations: an expectation of unit transverse residues on $n=5$
kissing cuts (planarity forces zeros instead), and the expectation
that quasi-linearity would first appear at $n=5$ or $6$ (it does not).
Third, everything here is exact, small, and reproducible: the
computations run in minutes on a laptop from the public repository, and
we would welcome correction or extension by practitioners --- in
particular a proof or counterexample for the all-$n$ conjecture, and
the composite-residue census.

\section*{Acknowledgments}
Computations and drafting were carried out in collaboration with an AI
research assistant (Claude, Anthropic) under the author's direction.
The debt to
refs.~\citep{ArkaniHamed:2013jha,ArkaniHamed:2013kca,ArkaniHamed:2010gh,
Kojima:2020gxs,ArkaniHamed:2017tmz,Dian:2024tvi} is evident throughout.

\bibliographystyle{unsrtnat}
\bibliography{refs}

\end{document}
